% GENERATED FILE. Edit canonical lessons, then run npm run generate:books.
% canonical-source-hash: fnv1a64:881f0f3ee5916045
% canonical-lessons: TA-C340-vilcciyatai, TA-C340-tura, TA-C340-natukatattu, TA-C340-porkalam, TA-C340-nakarikam

\chapter{Rise and Fall}
\label{ch:ta-rise-and-fall}
\hlchaptermodality{340}

\begin{chapteropening}
\textbf{By the end of this chapter:} \emph{I can say how kingdoms rose and fell.}

Complete the last lesson of chapter 340: I can say how kingdoms rose and fell.
\end{chapteropening}

\section[\texorpdfstring{\ta{வீழ்ச்சியடை} (vīḻcciyaṭai)}{வீழ்ச்சியடை (vīḻcciyaṭai)}]{\ta{வீழ்ச்சியடை} (vīḻcciyaṭai) — to decline; to fall}

\label{lesson:TA-C340-vilcciyatai}



\begin{quote}
Before the new one: say the Tamil for a pioneer, then the Tamil for a hero stone.
\end{quote}

\subsection*{You'll want to know: \ta{வீழ்ச்சியடை}}
\textbf{\ta{வீழ்ச்சியடை}} — \emph{vīḻcciyaṭai} — \textquotedblleft{}to decline; to fall\textquotedblright{}.

\begin{quote}
\ta{நூறு} \ta{ஆண்டுகளில்} \ta{அந்தப்} \ta{பேரரசு} \ta{வீழ்ச்சியடைந்தது}.
\end{quote}

\begin{quote}
\emph{nūṟu āṇṭukaḷil antap pēraracu vīḻcciyaṭaintatu.}
\end{quote}

\begin{quote}
In a hundred years that empire declined.
\end{quote}

Past (it …): \ta{வீழ்ச்சியடைந்தது}, \emph{vīḻcciyaṭaintatu}.

\ta{வீழ்ச்சி}, a fall + \ta{அடை}, to reach.

\begin{cousinweb}
\textbf{\ta{வீழ்ச்சியடை}} (\emph{vīḻcciyaṭai}) is a Tamil compound built from \ta{வீழ்ச்சி}, \textquotedblleft{}a fall\textquotedblright{}, from \ta{வீழ்}, \textquotedblleft{}to fall\textquotedblright{} (DEDR 5430), and \ta{அடை}, \textquotedblleft{}to reach\textquotedblright{} (DEDR 79).
\end{cousinweb}

\begin{grammarlens}[title={What you've built}]
One more word for telling what happened.
\end{grammarlens}

\subsection*{Your turn}
\begin{itemize}
\raggedright
  \item \emph{Say it:} \emph{vīḻcciyaṭai}
  \item \emph{Say it:} \emph{vīḻcciyaṭai}, once more
  \item \emph{Say it:} \emph{nūṟu āṇṭukaḷil antap pēraracu vīḻcciyaṭaintatu.} — In a hundred years that empire declined.
  \item \emph{Recall:} say the Tamil for a pioneer, then the Tamil for a hero stone, then say \emph{vīḻcciyaṭai} again
\end{itemize}

\begin{tcolorbox}[breakable,colback=teal!4,colframe=teal!35!black,title={Before you move on}]
What does \ta{வீழ்ச்சியடை} mean? (\textquotedblleft{}To decline; to fall\textquotedblright{}.) Say it once more.
\end{tcolorbox}
\section[\texorpdfstring{\ta{துற} (tuṟa)}{துற (tuṟa)}]{\ta{துற} (tuṟa) — to renounce; to give up}

\label{lesson:TA-C340-tura}



\begin{quote}
Before the new one: say the Tamil for a hero stone, then the Tamil for to decline.
\end{quote}

\subsection*{You'll want to know: \ta{துற}}
\textbf{\ta{துற}} — \emph{tuṟa} — \textquotedblleft{}to renounce; to give up\textquotedblright{}.

\begin{quote}
\ta{மன்னர்} \ta{அரியணையைத்} \ta{துறந்து} \ta{காட்டுக்குச்} \ta{சென்றார்}.
\end{quote}

\begin{quote}
\emph{maṉṉar ariyaṇaiyait tuṟantu kāṭṭukkuc ceṉṟār.}
\end{quote}

\begin{quote}
The king gave up the throne and went to the forest.
\end{quote}

Past (I …): \ta{துறந்தேன்}, \emph{tuṟantēṉ}.

What is given up takes -\ta{ஐ}; \ta{துறவி}, a monk, is one who has renounced.

\begin{cousinweb}
\textbf{\ta{துற}} (\emph{tuṟa}) is native Dravidian (DEDR 3365): to leave, relinquish, forsake; it takes its past in -\ta{ந்த}-, \ta{துறந்தேன்}.
\end{cousinweb}

\begin{grammarlens}[title={What you've built}]
One more word for telling what happened.
\end{grammarlens}

\subsection*{Your turn}
\begin{itemize}
\raggedright
  \item \emph{Say it:} \emph{tuṟa}
  \item \emph{Say it:} \emph{tuṟa}, once more
  \item \emph{Say it:} \emph{maṉṉar ariyaṇaiyait tuṟantu kāṭṭukkuc ceṉṟār.} — The king gave up the throne and went to the forest.
  \item \emph{Recall:} say the Tamil for a hero stone, then the Tamil for to decline, then say \emph{tuṟa} again
\end{itemize}

\begin{tcolorbox}[breakable,colback=teal!4,colframe=teal!35!black,title={Before you move on}]
What does \ta{துற} mean? (\textquotedblleft{}To renounce; to give up\textquotedblright{}.) Say it once more.
\end{tcolorbox}
\section[\texorpdfstring{\ta{நாடுகடத்து} (nāṭukaṭattu)}{நாடுகடத்து (nāṭukaṭattu)}]{\ta{நாடுகடத்து} (nāṭukaṭattu) — to exile; to deport}

\label{lesson:TA-C340-natukatattu}



\begin{quote}
Before the new one: say the Tamil for to decline, then the Tamil for to renounce.
\end{quote}

\subsection*{You'll want to know: \ta{நாடுகடத்து}}
\textbf{\ta{நாடுகடத்து}} — \emph{nāṭukaṭattu} — \textquotedblleft{}to exile; to deport\textquotedblright{}.

\begin{quote}
\ta{மன்னர்} \ta{தன்} \ta{தம்பியை} \ta{நாடுகடத்தினார்}.
\end{quote}

\begin{quote}
\emph{maṉṉar taṉ tampiyai nāṭukaṭattiṉār.}
\end{quote}

\begin{quote}
The king exiled his younger brother.
\end{quote}

Past (I …): \ta{நாடுகடத்தினேன்}, \emph{nāṭukaṭattiṉēṉ}.

\ta{நாடு}, country + \ta{கடத்து}, to send across.

\begin{cousinweb}
\textbf{\ta{நாடுகடத்து}} (\emph{nāṭukaṭattu}) is a Tamil compound built from \ta{நாடு}, \textquotedblleft{}country\textquotedblright{} (DEDR 3638), and \ta{கடத்து}, \textquotedblleft{}to make cross\textquotedblright{} (DEDR 1109).
\end{cousinweb}

\begin{grammarlens}[title={What you've built}]
One more word for telling what happened.
\end{grammarlens}

\subsection*{Your turn}
\begin{itemize}
\raggedright
  \item \emph{Say it:} \emph{nāṭukaṭattu}
  \item \emph{Say it:} \emph{nāṭukaṭattu}, once more
  \item \emph{Say it:} \emph{maṉṉar taṉ tampiyai nāṭukaṭattiṉār.} — The king exiled his younger brother.
  \item \emph{Recall:} say the Tamil for to decline, then the Tamil for to renounce, then say \emph{nāṭukaṭattu} again
\end{itemize}

\begin{tcolorbox}[breakable,colback=teal!4,colframe=teal!35!black,title={Before you move on}]
What does \ta{நாடுகடத்து} mean? (\textquotedblleft{}To exile; to deport\textquotedblright{}.) Say it once more.
\end{tcolorbox}
\section[\texorpdfstring{\ta{பொற்காலம்} (poṟkālam)}{பொற்காலம் (poṟkālam)}]{\ta{பொற்காலம்} (poṟkālam) — a golden age}

\label{lesson:TA-C340-porkalam}



\begin{quote}
Before the new one: say the Tamil for to renounce, then the Tamil for to exile.
\end{quote}

\subsection*{You'll want to know: \ta{பொற்காலம்}}
\textbf{\ta{பொற்காலம்}} — \emph{poṟkālam} — \textquotedblleft{}a golden age\textquotedblright{}.

\begin{quote}
\ta{அது} \ta{தமிழ்} \ta{இலக்கியத்தின்} \ta{பொற்காலம்}.
\end{quote}

\begin{quote}
\emph{atu tamiḻ ilakkiyattiṉ poṟkālam.}
\end{quote}

\begin{quote}
That was the golden age of Tamil literature.
\end{quote}

\ta{பொன்}, gold + \ta{காலம்}; \ta{ன்} turns to \ta{ற்} before \ta{க}.

\begin{cousinweb}
\textbf{\ta{பொற்காலம்}} (\emph{poṟkālam}) joins native \ta{பொன்}, \textquotedblleft{}gold\textquotedblright{} (DEDR 4570), with \ta{காலம்}, \textquotedblleft{}time\textquotedblright{}, from Sanskrit \emph{kāla} (CDIAL 3084).
\end{cousinweb}

\begin{grammarlens}[title={What you've built}]
One more word for telling what happened.
\end{grammarlens}

\subsection*{Your turn}
\begin{itemize}
\raggedright
  \item \emph{Say it:} \emph{poṟkālam}
  \item \emph{Say it:} \emph{poṟkālam}, once more
  \item \emph{Say it:} \emph{atu tamiḻ ilakkiyattiṉ poṟkālam.} — That was the golden age of Tamil literature.
  \item \emph{Recall:} say the Tamil for to renounce, then the Tamil for to exile, then say \emph{poṟkālam} again
\end{itemize}

\begin{tcolorbox}[breakable,colback=teal!4,colframe=teal!35!black,title={Before you move on}]
What does \ta{பொற்காலம்} mean? (\textquotedblleft{}A golden age\textquotedblright{}.) Say it once more.
\end{tcolorbox}
\section[\texorpdfstring{\ta{நாகரிகம்} (nākarikam)}{நாகரிகம் (nākarikam)}]{\ta{நாகரிகம்} (nākarikam) — a civilization; refinement}

\label{lesson:TA-C340-nakarikam}



\begin{quote}
Before the new one: say the Tamil for to exile, then the Tamil for a golden age.
\end{quote}

\subsection*{You'll want to know: \ta{நாகரிகம்}}
\textbf{\ta{நாகரிகம்}} — \emph{nākarikam} — \textquotedblleft{}a civilization; refinement\textquotedblright{}.

\begin{quote}
\ta{ஆற்றங்கரைகளில்} \ta{பழைய} \ta{நாகரிகங்கள்} \ta{தோன்றின}.
\end{quote}

\begin{quote}
\emph{āṟṟaṅkaraikaḷil paḻaiya nākarikaṅkaḷ tōṉṟiṉa.}
\end{quote}

\begin{quote}
Ancient civilizations arose on river banks.
\end{quote}

Also good manners: \ta{நாகரிகமாகப்} \ta{பேசு}, to speak politely.

\begin{cousinweb}
\textbf{\ta{நாகரிகம்}} (\emph{nākarikam}) is built on Sanskrit \emph{nāgarika}, \textquotedblleft{}of the town, a citizen\textquotedblright{} (CDIAL 7043): the ways of town life.
\end{cousinweb}

\begin{grammarlens}[title={What you've built}]
One more word for telling what happened.
\end{grammarlens}

\subsection*{Your turn}
\begin{itemize}
\raggedright
  \item \emph{Say it:} \emph{nākarikam}
  \item \emph{Say it:} \emph{nākarikam}, once more
  \item \emph{Say it:} \emph{āṟṟaṅkaraikaḷil paḻaiya nākarikaṅkaḷ tōṉṟiṉa.} — Ancient civilizations arose on river banks.
  \item \emph{Recall:} say the Tamil for to exile, then the Tamil for a golden age, then say \emph{nākarikam} again
\end{itemize}

\begin{tcolorbox}[breakable,colback=teal!4,colframe=teal!35!black,title={Before you move on}]
What does \ta{நாகரிகம்} mean? (\textquotedblleft{}A civilization; refinement\textquotedblright{}.) Say it once more.
\end{tcolorbox}
